% Supplement: additional numerical data.
% Source: draft/bie-enclosure/s-results.tex (clean version); same runs as
% \cref{sec:results}.
\section{Additional numerical data}\label{sm:results}

\subsection{Example 1}
All groups use $\ell_g=1$, $N_g=1$ and $\lambda_0=1/68$; the supremum
\eqref{eq:Cg} is attained at the stationary point for all groups.
\Cref{sm:tab-trace} lists the trace inputs. The curvature bounds hold on parameter
balls containing the binary64 and the decimal geometry; $\delta_g$ and $\eta_g$
come from a floating-point scan. \Cref{sm:tab-tails} lists the lead parts of the
defects. Without a center part, the lead rod groups equal their side values; the
center parts of the wall and seam groups are $1.966\times10^{-13}$ and
$2.35\times10^{-15}$. The mass bound uses disks of radius $0.165$ in the star, $0.168$
and $0.152$ in the two neighboring lead rods, and $0.037$--$0.070$ in the five lobes;
the center values are evaluated by $30$-point Gauss quadrature with error at most
$7.8\times10^{-8}$, and the three large disks give $0.543$, $0.392$, $0.347$, so
$\underline U^{\,2}\ge1.5273$. The tightest of the $472$ inertia tests has margin
$1.4\times10^{-4}$.

\begin{table}[tbhp]
\caption{Example 1: inputs of \cref{prop:trace} (displayed values rounded).}
\label{sm:tab-trace}
\centering\small
\begin{tabular}{lcccccc}
\toprule
Group & $\kappa_g$ & $\delta_g$ & $h_g$ & $\eta_g$ & $D_{0,g}$ & $D_{1,g}$\\
\midrule
center rod & $64.3$ & $0.015070$ & $0.969$ & $0.042232$ & $34.579$ & $0.20514$\\
left rods & $8.81$ & $0.094357$ & $0.831$ & $0.38312$ & $7.3292$ & $0.075821$\\
right rods & $10.53$ & $0.081458$ & $0.858$ & $0.31178$ & $8.0519$ & $0.080976$\\
walls, seams & $0$ & $0.2$ & $0$ & $1.1658$ & $5.4144$ & $0.061927$\\
\bottomrule
\end{tabular}
\end{table}

\begin{table}[tbhp]
\caption{Example 1: lead parts. The last column bounds
$(\sum_jb_{g,\pm}(\mathsf P_\pm^jc_\pm))^{1/2}$ including the tail.}
\label{sm:tab-tails}
\centering\small
\begin{tabular}{llcccc}
\toprule
Side & Group & $N_{\mathrm{cell}}$ & $\overline B_\pm$ & tail/finite (squared) & bound\\
\midrule
left & rods & $52$ & $2.72\times10^{5}$ & $3.0\times10^{-7}$ & $5.0811\times10^{-14}$\\
left & walls & $54$ & $1.32\times10^{6}$ & $1.7\times10^{-7}$ & $4.6633\times10^{-14}$\\
left & seams & $56$ & $1.91\times10^{6}$ & $8.5\times10^{-7}$ & $4.2586\times10^{-15}$\\
right & rods & $56$ & $5.78\times10^{6}$ & $3.9\times10^{-7}$ & $5.2316\times10^{-14}$\\
right & walls & $56$ & $2.51\times10^{6}$ & $3.4\times10^{-7}$ & $2.4254\times10^{-14}$\\
right & seams & $58$ & $3.62\times10^{6}$ & $9.3\times10^{-7}$ & $3.1306\times10^{-15}$\\
\bottomrule
\end{tabular}
\end{table}

\Cref{sm:tab-resolution} varies the field: each row is a complete cold-start run.
\Cref{sm:tab-eval} varies only the evaluation of the field of the third row.
Doubling all panels, $L=64$ or $192$ bits change every group bound by a relative
$3\times10^{-6}$ at most; the reliable bounds exceed a plain floating-point
quadrature of the residual polynomials by at most a relative $10^{-5}$. The five
candidates of \cref{sm:tab-resolution} differ from $1.558437894182088$ by $-2$ to
$+3$ ulps. The third row reproduces the main run bit for bit; its time differs
($919$ versus $827$\,s) because of machine load.

\begin{table}[tbhp]
\caption{Example 1, field resolution: nodes on the center curve (half as many
on each lead curve), $M$,
group bounds, $Q$, full width and time (upward rounded).}
\label{sm:tab-resolution}
\centering\scriptsize\setlength{\tabcolsep}{1.5pt}
\begin{tabular}{lccccccccc}
\toprule
Nodes & $M$ & center & left & right & walls & seams & $Q$ & width & s\\
\midrule
$256$ & $48$ & $2.2\times10^{-7}$ & $4.8\times10^{-14}$ & $9.7\times10^{-11}$ & $2.9\times10^{-13}$ & $1.1\times10^{-14}$ & $5.5\times10^{-7}$ & $4.28\times10^{-6}$ & $778$\\
$384$ & $48$ & $5.1\times10^{-10}$ & $4.9\times10^{-14}$ & $5.2\times10^{-14}$ & $2.0\times10^{-13}$ & $6.8\times10^{-15}$ & $1.3\times10^{-9}$ & $9.87\times10^{-9}$ & $764$\\
$512$ & $48$ & $1.3\times10^{-12}$ & $5.1\times10^{-14}$ & $5.3\times10^{-14}$ & $2.1\times10^{-13}$ & $5.8\times10^{-15}$ & $3.5\times10^{-12}$ & $2.738\times10^{-11}$ & $919$\\
$512$ & $40$ & $1.3\times10^{-12}$ & $5.2\times10^{-14}$ & $5.3\times10^{-14}$ & $2.0\times10^{-13}$ & $5.3\times10^{-15}$ & $3.5\times10^{-12}$ & $2.735\times10^{-11}$ & $812$\\
$512$ & $56$ & $1.3\times10^{-12}$ & $5.0\times10^{-14}$ & $5.2\times10^{-14}$ & $2.4\times10^{-13}$ & $1.0\times10^{-14}$ & $3.6\times10^{-12}$ & $2.766\times10^{-11}$ & $841$\\
\bottomrule
\end{tabular}
\end{table}

\begin{table}[tbhp]
\caption{Example 1, evaluation resolution for one field. Baseline: $192$ center
panels, $64$ per lead curve, $16$ per wall, $8$ per seam, $L=48$, $128$ bits. Times
cover the defect-to-assembly stages; the baseline time is the sum of these stages in
the complete run.}
\label{sm:tab-eval}
\centering\small
\begin{tabular}{lcccc}
\toprule
Setting & center $\overline\varepsilon_g$ & $Q$ & width & s\\
\midrule
baseline & $1.2615836\times10^{-12}$ & $3.4829355\times10^{-12}$ & $2.7378625\times10^{-11}$ & $598$\\
panels doubled & $1.2615800\times10^{-12}$ & $3.4829264\times10^{-12}$ & $2.7378554\times10^{-11}$ & $948$\\
$L=64$ & $1.2615836\times10^{-12}$ & $3.4829355\times10^{-12}$ & $2.7378625\times10^{-11}$ & $584$\\
$192$ bits & $1.2615836\times10^{-12}$ & $3.4829355\times10^{-12}$ & $2.7378625\times10^{-11}$ & $834$\\
\bottomrule
\end{tabular}
\end{table}

\begin{table}[tbhp]
\caption{Example 1: cost of one complete run ($512/256/256$) on an Apple M3 Pro
(11 cores, 18\,GiB), macOS 14, MATLAB R2023b, python-flint 0.9.0 (FLINT 3.6.0).
Wall-clock times include process start; peak memory is for the process tree.}
\label{sm:tab-cost}
\centering\small
\begin{tabular}{lrrr}
\toprule
Stage & s & parallelism & peak GiB\\
\midrule
gap classification and search & $175.6$ & MATLAB, 11 threads & $0.99$\\
density reconstruction and fit & $18.2$ & MATLAB, 11 threads & $1.17$\\
reliable defects (five groups) & $528.2$ & 8 processes & $1.82$\\
defect summary check & $1.1$ & 1 process & --\\
mass bound & $20.9$ & 1 process & $0.16$\\
trace constants and assembly & $1.1$ & 1 process & --\\
gap test & $80.5$ & 1 process & $0.21$\\
scheduling and output & $1.4$ & -- & --\\
\midrule
total & $827.1$ & & $1.84$\\
\bottomrule
\end{tabular}
\end{table}

\subsection{Example 2 and the gap-edge series}
\Cref{sm:tab-ex2} lists the groups for $L_0=1$. The collar checks use the actual
period $L_s$ of each side (smallest right margin $0.21$), and the seam integrals use
length $L_s$. The mass bound uses three disks in the empty center (centers $(0,0)$,
$(0,\pm0.33)$, radii $0.2$, $0.12$), where the field is a finite Rayleigh sum, so the
disk masses are computed from Bessel functions in ball arithmetic without
quadrature; $\underline U\ge0.645928$. In the gap test (window $k\in(1.46,2.37)$,
$944$ tests) the smallest margins are $2.92$ and $0.437$ at the lower ends and
$0.237$ and $1.20$ at the upper ends; the interval has margins $1.4469$ and $2.0383$.
The run takes $1058$\,s (search $453$\,s, defects $434$\,s, gap test $145$\,s) with
peak memory $2.24$\,GiB. With $384$ lead nodes the candidate has the same binary
representation; the lowest Rayleigh channel has $|\gamma_m|\ge1.82$.

\begin{table}[tbhp]
\caption{Example 2 ($L_0=1$, $\widehat k=1.8916961400527375$): groups and depths
$N_{\mathrm{cell}}$ (left, right).}
\label{sm:tab-ex2}
\centering\small
\begin{tabular}{lccccc}
\toprule
Group & $\overline C_g$ & $\overline\varepsilon_g$ & $\overline C_g\overline\varepsilon_g$ & left & right\\
\midrule
center rod & -- & $0$ & $0$ & -- & --\\
left rods & $1.135325$ & $2.0716\times10^{-14}$ & $2.3520\times10^{-14}$ & $42$ & --\\
right rods & $1.173284$ & $2.3041\times10^{-14}$ & $2.7033\times10^{-14}$ & -- & $52$\\
walls & $1.026040$ & $7.7580\times10^{-14}$ & $7.9600\times10^{-14}$ & $44$ & $51$\\
seams & $1.026040$ & $8.9003\times10^{-15}$ & $9.1320\times10^{-15}$ & $45$ & $52$\\
\midrule
$Q$ & & & $1.3929\times10^{-13}$ & &\\
\bottomrule
\end{tabular}
\end{table}

\Cref{sm:tab-float} is the floating-point continuation in $L_0$ (local Newton from the
previous point; smallest singular values below $2.2\times10^{-13}$, second smallest
$1.35$--$1.38$). The full intervals of the reliable runs are
\begin{align*}
  L_0=1.6:&\quad\mu\in[2.5009675413141,\,2.5009675413205],\\
  L_0=1.9:&\quad\mu\in[2.1912667743959,\,2.1912667744012],\\
  L_0=2:&\quad\mu\in[2.1186466555847,\,2.1186466555930],
\end{align*}
with margins to the lower window end $0.3693$, $0.0596$, $0.0161$ and to the upper
end $3.1159$, $3.4256$, $3.4982$. Increasing every $N_{\mathrm{cell}}$ by $10$ reduces
the widths by $0.06\%$, $1.3\%$, $0.86\%$. The predicted right rod depths
$\lceil\ln(\tau\underline V/(\overline B^2F_p\|c\|^2))/(2\ln r^{1/p})\rceil$ with
$F_p=\sum_{l<p}\|\mathsf P^l\|_2^2/(1-r^2)$ are $74$, $134$, $223$ versus the actual
$72$, $128$, $208$. For $L_0=2$, right side: $\|\mathsf P_+\|_2\le0.7953$,
$\|\mathsf P_+^2\|_2\le0.6219$, $F_1=2.7211$, $F_2=2.6616$. For both $p=1$ and $p=2$
the interval recomputed with $N_{\mathrm{cell}}+10$ meets the original one, and with
$p=2$ the next ten actual states plus the new tail stay below the old tail. The run
for $L_0=1.6$ stopped after the defects because of an input format error in the mass
stage; the remaining stages ran in a second process on the same field.

\begin{table}[tbhp]
\caption{Floating-point continuation in $L_0$ ($\varrho$ = spectral radius, $\delta_{\mathrm e}$ = distance of $\widehat k$ to the band edge).}
\label{sm:tab-float}
\centering\small
\begin{tabular}{cccccc}
\toprule
$L_0$ & $\widehat k$ & $\delta_{\mathrm e}$ & $\|\mathsf P_-\|_2$ & $\varrho(\mathsf P_+)$ & $\|\mathsf P_+\|_2$\\
\midrule
$1$ & $1.89170$ & $0.451$ & $0.2952$ & $0.3269$ & $0.3337$\\
$1.4$ & $1.66820$ & $0.227$ & $0.3323$ & $0.4092$ & $0.4173$\\
$1.6$ & $1.58144$ & $0.140$ & $0.3686$ & $0.4830$ & $0.4921$\\
$1.8$ & $1.51003$ & $0.0690$ & $0.4143$ & $0.5923$ & $0.6028$\\
$1.9$ & $1.48029$ & $0.0393$ & $0.4399$ & $0.6708$ & $0.6823$\\
$2$ & $1.45556$ & $0.0146$ & $0.4653$ & $0.7823$ & $0.7952$\\
\bottomrule
\end{tabular}
\end{table}
